\section*{Exercises}
\addcontentsline{toc}{section}{Exercises}

\subsection*{A. Theory}

\begin{exercise}
\exerciseid{1.A1}
Explain why the averaging rule in the chapter is linear. Your answer should
check both addition and scalar multiplication.
\end{exercise}

\begin{exercise}
\exerciseid{1.A2}
Give an example of a rule on grid states that is not linear. State which part
of linearity fails.
\end{exercise}

\subsection*{B. By Hand}

\begin{exercise}
\exerciseid{1.B1}
Let $u=(1,0,0,0)$ on a four-cell circular grid. Apply the averaging rule
$(Tu)_j=(u_{j-1}+u_{j+1})/2$ once by hand.
\end{exercise}

\subsection*{C. Python / NumPy}

\begin{exercise}
\exerciseid{1.C1}
Use basic \Python\ and \NumPy\ to implement the averaging rule for vectors of
length $N$. Test that your function satisfies
\texttt{T(u + v) == T(u) + T(v)} up to numerical tolerance.
\end{exercise}

\subsection*{D. Challenge}

\begin{exercise}
\exerciseid{1.D1}
Find a nonzero vector on a four-cell circular grid that is unchanged by the
averaging rule.
\end{exercise}
